% !Mode:: "TeX:UTF-8"
% !TEX program  = xelatex
%
% SlimeVR DIY 全身动作捕捉项目 —— 独立技术文档（单一源）
% 本文档独立描述本项目的目标、架构、关键技术与实现路线，
% 不依附于任何课程项目申报材料。后续所有设计修订、实现记录与评测结果均追加于此。
\documentclass[scheme=plain,UTF8,11pt,a4paper]{ctexart}

\usepackage[margin=1in]{geometry}
\usepackage{amsmath, amssymb, amsthm}
\usepackage{booktabs}
\usepackage{multirow}
\usepackage{makecell}
\usepackage{graphicx}
\usepackage{float}

% 架构图（单色/中性，不使用彩色）
\usepackage{tikz}
\usetikzlibrary{shapes.geometric, arrows.meta, positioning, calc, fit, backgrounds}
\tikzstyle{anchorbox} = [rectangle, draw=black, thick, rounded corners, fill=gray!12, minimum height=1.0cm, align=center, font=\small]
\tikzstyle{solverbox} = [rectangle, draw=black, thick, rounded corners, fill=white, minimum height=1.0cm, align=center, font=\small]
\tikzstyle{imubox}    = [rectangle, draw=black, thick, rounded corners, fill=gray!6,  minimum height=0.9cm, align=center, font=\small]
\tikzstyle{residbox}  = [rectangle, draw=black, thick, rounded corners, fill=white, minimum height=0.9cm, align=center, font=\small]
\tikzstyle{arrow}     = [-{Stealth}, thick, draw=black]
% 链接（隐藏颜色，保持单色）
\usepackage[hidelinks]{hyperref}

\title{\LARGE\bfseries 基于 SlimeVR 的低成本全身动作捕捉\\[0.3em]与服务器端神经网络误差纠正\\[0.6em]\large 系统设计与实现文档}

\date{\today}

\begin{document}

\maketitle

\begin{abstract}
本项目是一套自制的低成本全身动作捕捉系统：以 nRF52840 主控与 LSM6DSV 六轴 IMU 构建 6 个追踪节点，覆盖胸、髋、两大腿与两小腿，经 Wi-Fi 接入 SlimeVR-Server 解算骨架，最终驱动 VRChat 虚拟形象。系统的核心设计是在 PC 服务端引入一层轻量神经网络，作为“姿态先验 / 去噪器”，对多传感器姿态做联合误差纠正，弥补官方服务端仅依赖骨骼几何约束与最小二乘自动测比例（AutoBone）的不足。

由于所选 IMU 为六轴（无磁力计），偏航（yaw）只能由陀螺仪积分获得，其零偏与角随机游走会随时间累积，表现为全身缓慢转偏。对此，本设计采用“上半身绝对锚 + 下半身残差兜底”的分层抑制架构：上半身由 Virtual Desktop Estimation 生成的虚拟追踪器提供绝对、零漂移的姿态参考，下半身残留的局部小漂则由 SlimeVR 内置 Stay Aligned、本设计的神经网络残差校正器与骨架几何约束共同吸收。

校正器采用严格因果（无未来帧）的时空主干，单次前向并行输出姿态残差、传感器置信门控与足部触地判别三个结果，并经轻量物理仲裁层（加权 SLERP + ZUPT 零速锁脚）输出。训练走“纯离线数据驱动”范式：以 AMASS 光学动捕为纯净真值，注入 LSM6DSV 量级的陀螺零偏游走与噪声构造全合成配对样本，无需真人标注，也无需等待硬件到位；泛化能力由真实物理传感器数据集 DIP-IMU 独立评估。目标指标为单帧推理 $\le 2.0$\,ms（GPU FP16）、平均关节位置误差 MPJPE $\le 45$\,mm、偏航漂移消除率 $\ge 88\%$、足部触地判别 F1 $\ge 93\%$。端到端管线延迟控制在 20\,ms 以内。
\end{abstract}

\noindent\textbf{文档说明：}本文档是 SlimeVR DIY 全身动作捕捉项目的\textbf{主文档}，独立描述其目标、系统架构、硬件与固件、漂移抑制方案、校正器设计、方案定位与实施路线，\textbf{不依附于任何课程项目申报材料}。神经网络模型的训练数据集与实施规范另见同目录下的独立配套文档《SlimeVR 神经网络残差校正器 —— 模型训练、数据集与技术路线实施规范》（\texttt{model\_training\_specification.tex}）。除这两份文档外，不再另建独立文档；主文档层面的设计修订、实现记录与评测结果均追加于本文件。

\newpage
\tableofcontents
\newpage

% ==================== 1 项目概述 ====================
\section{项目概述}
\label{sec:overview}

本项目基于 SlimeVR 开源方案，自制一套以 nRF52840 + LSM6DSV 为核心的 6 点全身动作捕捉系统，用于 VRChat 虚拟形象驱动；核心设计是在 PC 服务端引入轻量神经网络，对多传感器姿态做联合误差纠正。

\textbf{动机。}商用全身追踪方案价格高昂，而 SlimeVR 开源硬件可将整套成本压到千元以内，显著降低虚拟形象驱动的门槛。

\textbf{创新点。}SlimeVR 官方服务端仅依赖骨骼几何约束、最小二乘自动测比例（AutoBone）与脚锁做纠偏：各 tracker 的绝对朝向由固件侧 VQF 融合给出，服务端信任该朝向并装配成骨架，因此既无跨 tracker 的统计融合，也不自动纠正偏航（yaw）漂移，单点失稳仅由相邻关节约束兜底。本设计在服务端增加一层神经网络姿态先验，用于抑制抖动、穿模与偏航漂移。

% ==================== 2 总体架构 ====================
\section{系统总体架构}
\label{sec:arch}

系统的数据链路为：6 个 IMU 追踪节点经 Wi-Fi 将姿态送入 SlimeVR-Server；Virtual Desktop Estimation 生成的上半身虚拟追踪器在 SteamVR 中以设备形式暴露，经官方 SlimeVR Feeder App 通过命名管道桥接进\emph{同一个} Server；两者在同一套骨架求解器中统一解算后，经 OpenVR 驱动或 OSC/VMC 输出至 VRChat。神经网络残差校正器以独立旁路进程运行，从 OSC/VMC 收取骨架、清洗后再回写，不改动 SlimeVR 代码。整体数据流如图~\ref{fig:pipeline} 所示。

\begin{figure}[H]
\centering
\begin{tikzpicture}[scale=0.82, every node/.style={transform shape}]
\node[imubox]    (imu)    at (0,0)     {6× IMU Tracker\\（nRF52840 + LSM6DSV）};
\node[imubox]    (vd)     at (7.2,0)   {VD Estimation 虚拟 Tracker\\（胸 / 上脊柱，绝对偏航）};
\node[solverbox] (server) at (3.6,-2.3){SlimeVR-Server 骨架求解器\\（统一解算 IMU + VD 设备）};
\node[solverbox] (out)    at (3.6,-4.2){OpenVR 驱动 / OSC·VMC 输出};
\node[residbox]  (nn)     at (9.0,-4.2){NN 残差校正\\旁路进程};
\node[solverbox] (vr)     at (3.6,-6.0){SteamVR $\to$ VRChat};

\draw[arrow] (imu) -- (server) node[midway, left]  {Wi-Fi};
\draw[arrow] (vd)  -- (server) node[midway, right] {Feeder App};
\draw[arrow] (server) -- (out);
\draw[arrow] (out) -- (vr);
\draw[arrow] (out) -- (nn) node[midway, above] {OSC/VMC};
\draw[arrow] (nn)  -- (vr) node[midway, right] {清洗回写};
\end{tikzpicture}
\caption{系统端到端数据流与神经网络旁路集成方式}
\label{fig:pipeline}
\end{figure}

% ==================== 3 硬件设计 ====================
\section{硬件设计}
\label{sec:hw}

6 个 tracker 分布于胸、髋、两大腿与两小腿，全部采用 6 轴 LSM6DSV（三轴加速度计 + 三轴陀螺仪，无磁力计），主控为 nRF52840，运行基于 Zephyr / nRF Connect SDK 的固件，经 Wi-Fi 将姿态数据回传至 PC 端 SlimeVR-Server。

每个追踪节点由 nRF52840 主控板、LSM6DSV IMU、电池与外壳构成；整套系统另需一个无线接收器用于汇聚 6 路数据，以及 DAPLink 调试器用于烧录与救砖（见第~\ref{sec:firmware} 节）。

% ==================== 4 固件与工具链 ====================
\section{固件与工具链}
\label{sec:firmware}

固件基于 nRF Connect SDK（Zephyr）构建，追踪器与接收器需使用\emph{不同}的 NCS 版本，因此各自建立独立的 west 工作区。该环节的主要风险与对策如下：

\begin{itemize}
    \item \textbf{IMU 识别待实机验证}：本项目所用器件在固件器件表中以 \textbf{LSM6DSV} 登记（I2C 地址 \texttt{0x70}），因此本文档全文统一称其为 LSM6DSV。但具体批次实际回报的器件 ID 仍需实机确认（已设为闸门 G1）：需通过实机 \texttt{info} 输出核对，若回报 \texttt{0x71} 则会落到 ISM330BX 驱动分支，需另行评估。
    \item \textbf{编译环境庞大}：NCS 工具链体积大，且追踪器与接收器版本不同，环境搭建易出错，需固定 commit 并单独记录。
    \item \textbf{刷写风险}：焊接与烧录顺序错误可能导致设备变砖；已准备 DAPLink 作为救砖手段。
    \item \textbf{算力边界}：nRF52840 无 NPU、算力不足，神经网络模型不应下沉到设备端，推理统一放在 PC 侧。
\end{itemize}

% ==================== 5 漂移与分层抑制 ====================
\section{偏航漂移问题与分层抑制架构}
\label{sec:drift}

\subsection{漂移机理}
\label{subsec:drift-mech}

本项目全部采用 6 轴 IMU（加速度计 + 陀螺仪，无磁力计），因此无法获得绝对偏航参考：加速度计只能依靠重力校正俯仰（pitch）与横滚（roll），而偏航（绕重力轴的旋转）只能由陀螺仪积分得到。陀螺零偏与角随机游走随时间累积（约 $0.1^{\circ}$/s 量级，即 $6^{\circ}$/min），表现为全身缓慢转偏。为此本设计引入“上半身绝对锚 + 下半身残差兜底”的分层抑制架构。

\subsection{组件与数据流}

上半身姿态（胸 / 肩 / 肘 / 腕）由 Virtual Desktop 的 Estimation 功能提供：它利用 Quest 3 头显自带的下视侧摄做 Inside-Out Body Tracking（IOBT），生成一组“虚拟 tracker”，其位姿为相机绝对测量、偏航零漂移，且零额外硬件。VD 虚拟 tracker 在 SteamVR 中以设备形式暴露，经由官方 SlimeVR Feeder App（随 SlimeVR 安装器附带、随 SteamVR 启动的 OpenVR 桥）通过命名管道喂入 SlimeVR-Server，与本项目的 6 个 IMU 在同一套骨架求解器中统一解算。社区工具 Virtual Desktop Body Tracking Configurator 可关闭 VD 的下半身虚拟 tracker（Generative Legs），仅保留上半身，把腿留给 IMU。

官方兼容文档明确 SlimeVR 与 Virtual Desktop Estimation “specifically for arms” 搭配使用。

\subsection{漂移抑制分层}

\begin{itemize}
    \item \textbf{第一层（全局锚，消除全局偏航漂移）}：VD Estimation 的胸 / 上脊柱（可含髋）作为绝对参考，骨架求解器自胸向下做前向运动学传播；胸一经固定，原本会贯穿全身的全局偏航漂移即被消除。
    \item \textbf{第二层（局部残差，吸收下半身小漂）}：仅余大腿 / 小腿的局部偏航小漂，幅度远小于全局漂，由以下三者共同吸收——(a) SlimeVR 内置 Stay Aligned（基于校准放松姿态的姿态先验，规则式梯度下降）；(b) 第~\ref{sec:nn} 节的神经网络残差校正器；(c) 骨架几何约束（骨长守恒、关节角限、地面约束、相邻 tracker 互推）。
\end{itemize}

\begin{figure}[H]
\centering
\begin{tikzpicture}[
    node distance=1.6cm and 3.2cm
]
\node[anchorbox] (vd) {VD Estimation 虚拟追踪器\\（胸/上脊柱，绝对偏航，零漂移）};
\node[solverbox] (solver) [below=of vd] {SlimeVR 骨架求解器\\（统一解算 IMU + VD 虚拟追踪器）};
\node[imubox]    (imu)    [left=of solver] {下半身 IMU trackers\\（大腿/小腿，局部偏航小漂）};
\node[residbox]  (resid)  [right=of solver] {残差吸收层（每 tick）：\\Stay Aligned（规则式）\\NN 残差校正器（本设计）\\骨架几何约束};
\node[solverbox] (vrc)    [below=of solver] {OSC / VMC 输出 $\to$ VRChat};

\draw[arrow] (vd) -- (solver) node[midway, left] {固定全局偏航};
\draw[arrow] (imu) -- (solver) node[midway, above] {局部小漂};
\draw[arrow] (solver) -- (vrc) node[midway, left] {骨架};
\draw[<->, arrow] (resid) -- (solver) node[midway, above] {每 tick 残差吸收};
\end{tikzpicture}
\caption{分层漂移抑制架构：上半身绝对锚消除全局偏航漂移，下半身局部小漂由残差吸收层兜底}
\label{fig:layered_arch}
\end{figure}

\subsection{Stay Aligned 与神经网络校正器的互补}

Stay Aligned 是规则式姿态先验：每 tick 将 tracker 偏航朝“校准时的放松姿态”做梯度下降微调，它读取的是玩家自身姿态与放松姿态先验，\textbf{不读取头显偏航}，因此本身并非绝对参考；其盲区在于非放松姿态（侧躺、大幅跳舞、前倾）时放松姿态不成立而失效。

神经网络残差校正器从 AMASS 等大规模动捕数据中学到“人类在各活动下的合理姿态分布”，可覆盖上述盲区。二者同属姿态先验类，神经网络是 Stay Aligned 的\emph{补强}而非替代；代价是神经网络可能产生“幻觉”（将姿态掰到合理但错误的位置），因此必须采用残差叠加 + 置信门控的形式（见~\ref{subsec:discipline}）。

% ==================== 6 神经网络残差校正器 ====================
\section{服务器端神经网络残差校正器}
\label{sec:nn}

\subsection{问题定位}

SlimeVR 官方服务端仅用骨骼几何约束、最小二乘自动测比例（AutoBone）与脚锁做纠偏，各 tracker 的绝对朝向由固件侧 VQF 融合给出，服务端信任该朝向并装配成骨架。其局限是：无跨 tracker 的统计融合、不自动纠正偏航（yaw）漂移、单点失稳仅由相邻关节约束兜底。

此外，纯开链几何解算缺少足部触地约束，静止站立时下肢易出现“滑步 / 溜冰”现象。因此本设计在 PC 服务端增加一层轻量神经网络，作为“姿态先验 / 去噪器”，同时输出姿态残差、传感器置信门控与足部触地判别，对多传感器姿态做联合误差纠正。

\subsection{输入与输出}

\textbf{输入：}滑动窗口特征 $\mathbf{X}_{\text{in}} \in \mathbb{R}^{N \times 135}$，由两路构成——(a) 上半身绝对锚：Quest 3 / VD Estimation 提供的胸、手肘绝对位姿（零漂移）；(b) 下半身 6 点 LSM6DSV 的角速度与加速度。两路数据先经时间戳 Slerp 环形插值缓冲池（50\,ms）对齐，以消除视觉传输（$\sim$20\,ms）与 2.4G 射频（$\sim$5\,ms）之间的相位差。

\textbf{输出（多任务头，单次前向并行给出）：}
\begin{itemize}
    \item \textbf{姿态残差头} $\Delta q$：各下半身关节的局部四元数残差，单帧角度限幅 $\le 8^{\circ}$；
    \item \textbf{可靠性门控头} $W$：传感器置信度，低置信时退避回原值；
    \item \textbf{触地判别头} $C$：左右脚触地概率 $P_L, P_R \in [0,1]$，供零速修正（ZUPT）与锁脚使用。
\end{itemize}

最终装配由轻量物理与运动学仲裁层完成：按置信度做加权 SLERP 叠加，并在双足同时触地时进行闭链支撑仲裁，避免单脚浮空或陷入地面。

\subsection{网络结构}

主干采用\textbf{因果时空网络}：8 路独立\emph{单向因果} LSTM（$256$ 维，严格不引入未来帧，保证在线推理零额外时延）+ Spatial Transformer 编码器（跨关节解剖关联注意力）。在此之上并列三个多任务头（残差 $\Delta q$ / 门控 $W$ / 触地 $C$），单次前向完成全部解算，避免运行时级联多个异构模型。

为避免重复造轮子，本项目复用 4 个开源仓库（已拉取至 \texttt{external\_repos/}），职责划分如表~\ref{tab:repos} 所示。

\begin{table}[H]
\centering
\small
\caption{核心参考仓库与本项目的技术映射}
\label{tab:repos}
\begin{tabular}{p{2.6cm}p{2.4cm}p{8.4cm}}
\toprule
\textbf{仓库} & \textbf{出处} & \textbf{在本项目中的复用职责} \\
\midrule
HMD-Poser     & CVPR 2024   & \textbf{主干网基石}：8 路身体流单向因果 LSTM 与 Spatial Transformer，实现零未来帧延迟的时空特征提取与多传感器可靠性门控 \\
PIP           & CVPR 2022   & \textbf{物理接触判决}：移植足部触地判别头（ContactNet）与 ZUPT 零速锁脚逻辑，消除下肢“溜冰感” \\
AvatarPoser   & ECCV 2022   & \textbf{数据工程脚手架}：复用 AMASS 序列切片、滑动窗口构造、SMPL-H 骨骼映射与归一化脚本 \\
SparsePoser   & SIGGRAPH 2023 & \textbf{解剖学流形参考}：借鉴关节极限约束（Joint Limit Loss）作为防止肢体非生理反曲的正则项 \\
\bottomrule
\end{tabular}
\end{table}

\textbf{备选路线}（供后期消融对比）：
\begin{itemize}
    \item \textbf{备选 A（因果膨胀卷积 TCN）}：将主干换为 TCN，推理无循环依赖，便于 ONNX Runtime 做 INT8 定点量化与算子融合，CPU 推理可压至 $0.8$\,ms 以内；
    \item \textbf{备选 B（姿态流形去噪后处理）}：在残差输出后串联轻量自编码器潜空间投影，关节角超限时投影回自然姿态流形；为防止“AI 幻觉吃动作”，须配合大残差自动退避机制。
\end{itemize}

模型总参数量控制在 $1.2$\,M 以内，导出 ONNX 权重不超过 $8$\,MB。

\subsection{训练数据与目标函数}

训练与验证全部离线完成，不依赖硬件进度。

\textbf{训练集（AMASS）}：以 MPI 发布的 AMASS 光学动捕为纯净真值（SMPL-H 关节旋转矩阵与全局平移）。精选三个子集即可覆盖日常动作，总量约 2.9\,GB，无需下载全量：\texttt{CMU}（约 1.5\,GB，走 / 跑 / 坐 / 转身 / 下蹲 / 跳跃）、\texttt{HDM05}（约 800\,MB，规范化高精度步态）、\texttt{BMLrub}（约 600\,MB，连续步态循环与原地踏步，是触地检测的最佳样本）。正向运动学计算需预先获取 SMPL-H 人体模型参数包。

\textbf{评估集（DIP-IMU）}：MPI 开源的 DIP-IMU（约 1.5\,GB）记录了 10 名被试佩戴 17 颗 Xsens 工业级 IMU 时的 Vicon 光学真值，可在硬件完成前就用\emph{真实}物理传感器数据检验模型对温漂、抖动与软组织晃动的泛化能力。

\textbf{合成加噪}：对每个下半身关节注入偏航随机游走
\[
\psi_t = \psi_{t-1} + b_{\text{gyro}}\Delta t + \sigma_{\text{arw}}\sqrt{\Delta t}\,\epsilon_t,\qquad \epsilon_t \sim \mathcal{N}(0,1)
\]
并叠加加速度扰动与高斯白噪声；同时由脚踝速度与离地高度生成触地真值（$v_{\text{foot}} < 0.1$\,m/s 且 $z_{\text{foot}} < 0.03$\,m 时置 1）。监督目标为残差 $\Delta q^{*} = q^{*} \otimes (q^{\text{corrupt}})^{-1}$，构成 corrupted→clean 配对，无需真人标注。

\textbf{损失与优化}：
\[
\mathcal{L} = \text{SmoothL1}(\Delta q, \Delta q^{*}) + \lambda_1\,\text{BCE}(P_{\text{contact}}, C^{*}) + \lambda_2\,\mathcal{L}_{\text{joint\_limit}}
\]
取 $\lambda_1 = 0.5$、$\lambda_2 = 0.1$；使用 AdamW 优化器（初始学习率 $10^{-3}$，Cosine 退火衰减），Batch Size 为 256，在消费级 GPU 上训练约 80 个 Epoch。

\subsection{与 SlimeVR 的集成方式}

\textbf{松耦合（推荐）：}SlimeVR 经 OSC / VMC 输出骨架 → 独立神经网络进程消费 → 清洗后回写 OSC / VMC 给 VRChat，不改动 SlimeVR 代码。工程上实现为一个轻量 Python 中继客户端：监听 SlimeVR Server 的 OSC 广播端口（6969 / 39539），用 ONNX Runtime 做实时推理，再将滤波与锁脚后的骨骼四元数重定向发送至 VRChat（端口 9000），即插即用、非侵入。

\textbf{紧耦合：}fork SlimeVR Server，在骨架计算后插入推理；需经 JNI / 子进程 / ONNX Runtime 桥接 Java 与模型。

\subsection{延迟与部署约束}

VR 要求端到端管线 $\le 20$\,ms（从 IMU 打包到 VRChat 收到骨骼）。其中神经网络单帧前向推理预算为 $\le 2.0$\,ms（GPU FP16），CPU 上的 ONNX 路径放宽至 $\le 4.5$\,ms；轻量物理与运动学仲裁层耗时控制在 $0.3$\,ms 以内。训练收敛后导出为静态计算图（ONNX，opset 14）并转换为 FP16，由 ONNX Runtime / TensorRT 执行。

部署边界明确：nRF52840 无 NPU、算力不足，模型\textbf{不}下沉到设备端，推理统一在 PC 侧完成。

\subsection{验证方案}

评测分两步：先在\emph{纯离线}条件下用 DIP-IMU 评估集跑量化指标（不连接任何硬件），达标后再上真机做稳定性与主观对比（抖动均方根、穿模率、滑步距离、reset 频率）。量化验收门槛如表~\ref{tab:kpi} 所示。

\begin{table}[H]
\centering
\small
\caption{校正器量化验收指标（完整训练与评测规范见配套文档）}
\label{tab:kpi}
\begin{tabular}{p{2.6cm}p{5.2cm}p{5.2cm}}
\toprule
\textbf{指标类别} & \textbf{关键参数项} & \textbf{验收门槛} \\
\midrule
\multirow{2}{*}{实时性}
    & 单帧前向推理时延 & $\le 2.0$\,ms（GPU FP16）/ $\le 4.5$\,ms（CPU ONNX） \\
    & 端到端全链路时延 & $\le 20.0$\,ms \\
\midrule
\multirow{3}{*}{姿态精度}
    & 平均关节位置误差 MPJPE & $\le 45.0$\,mm \\
    & 关节角度 RMSE & $\le 7.5^{\circ}$ \\
    & 偏航漂移消除率 & $\ge 88\%$（注入 $10^{\circ}$/min 漂移） \\
\midrule
\multirow{2}{*}{接触与稳定}
    & 足部触地判别 F1 & $\ge 93.0\%$ \\
    & 足部滑步距离削减 & $\ge 80\%$（相比官方开环） \\
\midrule
\multirow{2}{*}{轻量化}
    & 模型参数量 & $\le 1.2$\,M \\
    & ONNX 权重体积 & $\le 8.0$\,MB \\
\bottomrule
\end{tabular}
\end{table}

\subsection{实现纪律与评测}
\label{subsec:discipline}

为使神经网络稳妥，须遵守以下纪律：(1) 仅输出小幅度残差四元数 $\Delta q$ 并叠加回原朝向，绝不整体替换骨架；(2) 加置信门控，低置信时不修正、退回原值；(3) 训练数据全合成——以 AMASS 干净全身骨架为监督，经简化前向运动学重建后，在下半身关节注入 LSM6DSV 量级的偏航随机游走漂移（陀螺零偏 / 角随机游走）与 IMU 噪声，构造 corrupted→clean 配对回归残差，无需真人标注；(4) 结构以因果（单向）LSTM 主干起步，备选 A 的因果膨胀 TCN 用于消融对比；(5) 评测先量化后上线——在合成漂移上测均方误差与偏航误差恢复率（如注入 $10^{\circ}$/min 漂，看网络收敛回多少），达标后再上真机对比 reset 频率、穿模率与滑步距离；(6) 残差输出必须限幅（单帧角度 $\le 8^{\circ}$），并在出现异常大残差时自动退避，防止模型把姿态掰到“合理但错误”的位置。整体评测不依赖“感觉更稳”的主观判断。
% ==================== 7 对比与定位 ====================
\section{现有方案对比与项目定位}
\label{sec:compare}

经检索，SlimeVR 生态内\textbf{不存在}专门的“神经网络残差校正器”开源项目，该方向属生态空白；官方现有漂移缓解手段为 Stay Aligned（规则式）+ 偏航重置 + 可选磁力计。学术界有面向单 IMU 的 AHRS 神经网络基线（RoNIN、TLIO），但其未与 SlimeVR 集成；视觉 ML 全身追踪则提供另一条绝对参考路线。本设计的定位如表~\ref{tab:compare} 所示。

\begin{table}[H]
\centering
\small
\caption{现有方案与本设计的对比定位}
\label{tab:compare}
\begin{tabular}{p{3.0cm}p{2.6cm}p{2.0cm}p{4.0cm}}
\toprule
\textbf{方案} & \textbf{类型} & \textbf{是否绝对参考} & \textbf{局限 / 本项目定位} \\
\midrule
Stay Aligned & 规则式姿态先验 & 否（相对校准放松姿态） & 非放松姿态（侧躺/大幅舞蹈）失效；\textbf{保留为基线} \\
RoNIN / TLIO & 学术 AHRS 神经网络（单 IMU） & 否 & 单 IMU、无全身约束、需真机标定；\textbf{借鉴其网络思想} \\
相机 ML 全身追踪（KinectAIFusionFBT / Fever / Mediapipe-VR / ToucanTrack） & 视觉绝对参考 & 是 & 需额外摄像头/算力、遮挡敏感；\textbf{替代路线，本设计不采用} \\
\textbf{本设计}（VD 锚 + NN 残差校正） & 分层（绝对锚 + 统计残差） & 上半身绝对（VD）、下半身残差 & NN 幻觉需置信门控；\textbf{主路线} \\
\bottomrule
\end{tabular}
\end{table}

% ==================== 8 实施路线 ====================
\section{实施路线}
\label{sec:roadmap}

本项目按以下里程碑推进：

\begin{itemize}
    \item \textbf{阶段 1}：搭建开发环境，跑通 SlimeVR 软件四链（Server / GUI / OpenVR 驱动）。
    \item \textbf{阶段 2}：完成 6 点 tracker 的焊接组装与整机装配。
    \item \textbf{阶段 3}：烧录固件，实机完成闸门 G1 验证（IMU 识别）。
    \item \textbf{阶段 4}：设计并训练服务器端神经网络误差纠正模型。
    \item \textbf{阶段 5}：模型与 SlimeVR 集成，进行 VRChat 实机测试与调参。
    \item \textbf{阶段 6}：整理技术文档与成果归档。
\end{itemize}

\end{document}
